\Subhead{2.10.3}{Интенсивные акустические пучки. Параметрические излучатели звука}

В интенсивном пучке одновременно действуют дифракция, нелинейное укручение и поглощение. Для акустического давления $p$, сопутствующего времени $\tau=t-z/c_0$ одна из форм уравнения Хохлова–Заболотской–Кузнецова:
\begin{equation}
 \partial_z\partial_\tau p=\frac{c_0}{2}\Delta_\perp p+
 \frac{b}{2\rho_0c_0^3}\partial_\tau^3p+
 \frac{\beta}{2\rho_0c_0^3}\partial_\tau^2(p^2).
\end{equation}
$b$ — эффективная вязкость с теплопроводным вкладом. Уравнение предполагает малые углы, слабую нелинейность за период и медленное изменение по $z$; интегрирование по $\tau$ возвращает дифракционный, диссипативный и нелинейный члены.

Две близкие мощные частоты $\omega_1,\omega_2$ создают квадратичный источник на $\Omega=|\omega_1-\omega_2|$: произведение $p_1p_2$ содержит $\cos[(\omega_1-\omega_2)t]$. Длинная область их перекрытия становится распределённой антенной разностной частоты. Она может дать узкий низкочастотный пучок, хотя обычный компактный излучатель той же частоты имел бы широкую диаграмму. Накачки сильнее поглощаются, разностная волна распространяется дальше. Недостатки — малая эффективность и потребность в большой мощности; диаграмма определяется геометрией взаимодействия и поглощением, а не только длиной волны разностной частоты.
